\documentclass[10pt,a4paper]{article}

% ============================================================
% PACKAGES
% ============================================================
\usepackage[top=0.5in, bottom=0.45in, left=0.7in, right=0.7in]{geometry}
\usepackage{titlesec}
\usepackage{enumitem}
\usepackage{hyperref}
\usepackage{xcolor}
\usepackage[T1]{fontenc}
\usepackage{charter}
\usepackage{parskip}
\usepackage{needspace}
\usepackage{ragged2e}

% ============================================================
% COLORS (match raadh.me)
% ============================================================
\definecolor{heading}{HTML}{0B4F5C}
\definecolor{link}{HTML}{0B4F5C}
\definecolor{body}{HTML}{1A1A1A}
\definecolor{muted}{HTML}{5A6470}

% ============================================================
% HYPERLINKS
% ============================================================
\hypersetup{
    colorlinks=true,
    linkcolor=link,
    urlcolor=link,
    pdfauthor={Rokib Al Dhin Raadh},
    pdftitle={Rokib Al Dhin Raadh - Research Statement (a16z Frontier)},
    pdfsubject={AI researcher; founder and CEO of OXIEDO; inventor of ORMAS; DeepMath 2026; NeurIPS 2026 Workshop reviewer; Cohere Labs invited talk},
    pdfkeywords={Rokib Al Dhin Raadh, machine learning research, optimisation, training dynamics, evaluation, safety, publications, PyTorch, LLMs, ORMAS, DeepMath 2026},
}

% ============================================================
% SECTION STYLE
% ============================================================
\titleformat{\section}
  {\normalsize\bfseries\scshape\color{heading}}
  {}
  {0em}
  {}
  [\vspace{-1pt}{\color{heading}\titlerule[0.4pt]}]
\titlespacing{\section}{0pt}{9pt}{5pt}

\pagestyle{empty}

% ============================================================
% CV ENTRY: date in the left gutter, body indented
% ============================================================
% A list (not a minipage) so long entries can break across pages.
\newcommand{\cvitem}[3]{%
  \begin{list}{}{%
    \setlength{\leftmargin}{2.35cm}%
    \setlength{\labelwidth}{2.1cm}%
    \setlength{\labelsep}{0.25cm}%
    \setlength{\itemindent}{0pt}%
    \setlength{\listparindent}{0pt}%
    \setlength{\rightmargin}{0pt}%
    \setlength{\topsep}{0pt}%
    \setlength{\partopsep}{0pt}%
    \setlength{\parsep}{0pt}%
    \setlength{\itemsep}{0pt}%
  }%
  \item[{\makebox[2.1cm][l]{\footnotesize\color{muted}#1}}]%
  #2
  #3
  \end{list}%
  \vspace{1.5pt}
}

% Tight bullet list
\newlist{bullets}{itemize}{1}
\setlist[bullets]{leftmargin=1em, itemsep=0pt, parsep=0pt, topsep=1pt, label=\textbullet}

% Numbered publication list
\newlist{pubs}{enumerate}{1}
\setlist[pubs]{leftmargin=2.35cm, labelwidth=2.1cm, labelsep=0.25cm, labelindent=0pt, align=left, itemsep=2pt, parsep=0pt, topsep=1pt, label={\footnotesize\color{muted}[\arabic*]}}

\setlength{\emergencystretch}{2.5em}
\hyphenpenalty=10000
\exhyphenpenalty=50
\begin{document}
\RaggedRight
\color{body}
\begin{center}
  {\Large\bfseries\color{heading} Research Statement}\\[2pt]
  {\small Rokib Al Dhin Raadh \enspace$\cdot$\enspace \href{https://www.raadh.me}{raadh.me} \enspace$\cdot$\enspace \href{mailto:raadh@oxiedo.com}{raadh@oxiedo.com} \enspace$\cdot$\enspace a16z Frontier Research Intern, Summer 2027}
\end{center}

\section*{The question}
Can we know, while a model is training, that training is still sound, and prove it rather than guess it?

\section*{Why it is open}
Almost every convergence guarantee in deep learning assumes each update is a gradient step on a fixed loss. Real training breaks that assumption constantly. Gradient clipping, optimizer momentum, learning-rate schedules, data mixing, RL fine-tuning against a shifting reward, and self-repair all produce updates that are not gradients of anything. We train frontier models on mechanisms our theory does not cover, and we find failures afterwards: a loss spike that wastes a week of compute, a dead or saturated component, a reward hack that only shows up in evaluation.

\section*{Why it is hard}
The loss landscape is non-convex, so global guarantees are mostly out of reach. Non-gradient updates remove the usual tool, the loss itself as a Lyapunov function. And any check has to run inside a training job without slowing it down or drowning engineers in false alarms.

\section*{Why it matters}
Training is the most expensive and least observable step in building AI. A certificate that says ``this run is stable as long as these quantities stay bounded'', or ``these units can never recover'', turns monitoring from dashboards into decisions. It also matters for safety and regulation: an auditable record of what changed during training is the basis for trusting a model.

\section*{What I have done}
I have completed one full turn of this loop, alone.
\begin{bullets}
  \item \textbf{Theory.} In my DeepMath 2026 paper (sole author, double-blind) I treated self-repair as a bounded disturbance and proved local practical input-to-state stability, a control-theory guarantee that survives non-gradient updates. The same analysis gave a hard negative result: bounded per-channel repairs cannot revive a dead BatchNorm--ReLU channel. \href{https://openreview.net/forum?id=kowltBZEIv}{OpenReview}
  \item \textbf{Tool.} I turned that result into \href{https://github.com/Oxiedo/ormas-monitor}{ormas-monitor} (Apache-2.0, PyPI). On pretrained models it found 292 channels in VGG-19-BN that output exactly zero on all 10,000 ImageNetV2 photos; removing them cut compute by 18.6\% and changed no predictions, while removing 292 random live channels changed 499 of 500. \href{https://www.raadh.me/dead-channels}{Write-up}
  \item \textbf{Evidence.} 383 controlled, seeded experiments on one RTX~3090, all reproducible from code. I presented the work as an invited talk at Cohere Labs and published corrections to my own results there.
\end{bullets}

\section*{What I would do at a16z Frontier}
My results are on CNNs; the open question lives in transformers and RL post-training. In the internship I would (1)~derive stability conditions for common non-gradient updates such as clipping and RL fine-tuning, (2)~turn each condition into a cheap runtime check, and (3)~test those checks on real training runs, including failed ones, to measure how early they warn and how often they are wrong. The output would be a paper, an open-source release and a public write-up. a16z Frontier offers what I lack: researchers who bridge theory and practice, compute, and portfolio companies whose training runs are exactly the data this research needs.

\end{document}
